\documentclass[aspectratio=169,11pt]{beamer}
\usepackage[spanish,es-noshorthands]{babel}
\usepackage[utf8]{inputenc}
\usepackage[T1]{fontenc}
\usepackage{amsmath,amssymb}
\usepackage{hyperref}
\usetheme{Madrid}
\usecolortheme{dolphin}
\definecolor{crono}{RGB}{49,46,129}
\definecolor{acento}{RGB}{79,70,229}
\setbeamercolor{structure}{fg=crono}
\setbeamercolor{frametitle}{bg=crono,fg=white}
\setbeamercolor{title}{fg=white,bg=crono}
\setbeamertemplate{navigation symbols}{}
\setbeamertemplate{footline}[frame number]
\hypersetup{colorlinks=true,urlcolor=acento}
\title{CronoGraph}
\subtitle{Optimización de calendarios de parciales mediante grafos}
\author{Daniel Alejandro Rincón Navarro \and Matías Guarín Rueda \and Manuel David Polo Salgado \and Andrés David Issa Camargo \and Cristian Camilo Fuentes Sánchez}
\institute{Universidad del Norte \\ Matemáticas Discretas}
\date{}
\begin{document}

\begin{frame}
  \titlepage
  \centering\scriptsize Repositorio: \url{https://github.com/WosDan/CronoGraph}
\end{frame}

\begin{frame}{El problema y por qué lo elegimos}
  \begin{block}{La situación}
    Un estudiante puede tener varios parciales próximos, aunque no coincidan a la misma hora.
  \end{block}
  \vspace{0.3cm}
  \begin{itemize}
    \item Coordinar fechas manualmente se vuelve difícil cuando muchos estudiantes comparten materias.
    \item Elegimos el problema por su relación con la carga académica y porque se puede modelar con grafos.
    \item Objetivo: proponer calendarios con menor concentración, respetando la disponibilidad registrada.
  \end{itemize}
\end{frame}

\begin{frame}{Datos y representación con grafos}
  \begin{columns}[T,onlytextwidth]
    \begin{column}{0.52\textwidth}
      \textbf{Datos de demostración}\par\medskip
      120 estudiantes ficticios\\ 15 asignaturas y 12 docentes\\ 20 parciales con franjas autorizadas
      \medskip\par Cada matrícula conecta un estudiante con sus parciales.
    \end{column}
    \begin{column}{0.44\textwidth}
      \textbf{Dos grafos}\par\medskip
      1. Bipartito: estudiantes y parciales.\par\medskip
      2. Conflictos: parciales conectados por estudiantes compartidos.
      \[A=B^{\mathsf T}B\]
      \footnotesize $A_{ij}$ cuenta los estudiantes que presentan ambos parciales.
    \end{column}
  \end{columns}
\end{frame}

\begin{frame}{Cómo medimos la concentración}
  \begin{columns}[T,onlytextwidth]
    \begin{column}{0.48\textwidth}
      \textbf{Peso académico}
      \[W_i=C_i p_i\left(1+\frac{r_i}{2}\right)\]
      \footnotesize Créditos $C_i$, porcentaje del parcial $p_i$ e índice de retiro $r_i$.
    \end{column}
    \begin{column}{0.48\textwidth}
      \textbf{Proximidad}
      \[P(\Delta)=\max(0,4-\Delta)\]
      \footnotesize Se penalizan parciales separados por 0 a 3 días calendario.
    \end{column}
  \end{columns}
  \vspace{0.15cm}
  \[\min F=\sum_{i<j}A_{ij}P(\Delta_{ij})\frac{W_i+W_j}{2}\]
  \centering\footnotesize Los pesos son decisiones del modelo; no miden directamente el estrés.
\end{frame}

\begin{frame}{Cómo construimos una propuesta}
  \begin{enumerate}
    \item Partimos de un calendario y de las franjas autorizadas para cada parcial.
    \item Nuestra búsqueda prueba cambios factibles y conserva mejoras; usa varios puntos de partida.
    \item Las centralidades ayudan a ordenar algunos puntos de partida y a desempatar movimientos.
    \item Opcionalmente, CP-SAT intenta mejorar o verificar la propuesta.
  \end{enumerate}
  \vspace{0.2cm}
  \begin{block}{Restricciones}
    Franjas autorizadas; máximo 14 días de cambio; sin solapamientos de estudiantes, docentes o aulas.
  \end{block}
\end{frame}

\begin{frame}{Demostración del simulador}
  \begin{enumerate}
    \item Abrir \href{https://crono-graph.vercel.app/}{crono-graph.vercel.app} y revisar el calendario inicial.
    \item Ejecutar la optimización y comparar el horario propuesto.
    \item Seleccionar un parcial para revisar su franja autorizada.
    \item Abrir \textbf{Estudiantes} y comparar los pares de parciales cercanos.
  \end{enumerate}
  \vspace{0.25cm}
  \footnotesize La demostración utiliza datos ficticios. Cada resultado corresponde a la ejecución mostrada.
\end{frame}

\begin{frame}{Limitaciones y trabajo futuro}
  \begin{itemize}
    \item La búsqueda encuentra soluciones factibles, pero no garantiza el mínimo global.
    \item CP-SAT certifica el óptimo de su objetivo redondeado solo si alcanza \texttt{OPTIMAL}.
    \item Faltan datos reales y validación de los pesos y del avance de los temas de clase.
    \item Comparar algoritmos bajo las mismas condiciones y verificar instancias pequeñas por enumeración.
  \end{itemize}
\end{frame}

\begin{frame}{Conclusiones}
  \begin{itemize}
    \item Los grafos conectan las matrículas con los conflictos entre parciales.
    \item El simulador compara calendarios respetando las restricciones registradas.
    \item Si el calendario inicial es factible, nuestra búsqueda conserva uno de costo igual o menor.
    \item El siguiente paso es contrastar métodos y validar los criterios con datos representativos.
  \end{itemize}
  \vfill
  \small Repositorio: \url{https://github.com/WosDan/CronoGraph}\\
  Aplicación: \url{https://crono-graph.vercel.app/}
\end{frame}

\end{document}
